\documentclass[12pt]{article}
\usepackage[utf8]{inputenc}
\usepackage[T1]{fontenc}
\usepackage[margin=1in]{geometry}
\usepackage{setspace}
\usepackage{xcolor}
\usepackage[colorlinks=true, urlcolor=blue!60!black, linkcolor=blue!60!black]{hyperref}
\setlength{\parindent}{0em}
\setlength{\parskip}{0.6em}
\renewcommand{\baselinestretch}{1.0}
\begin{document}
\begin{center}
{\LARGE \textit{Pike v. Bruce Church, Inc.}, 397 U.S. 137 (1970)\par}
\vspace{0.5em}
{\large Case Brief --- Two-Page Edition\par}
\end{center}

\section{Memory Jogger}
A unanimous Court struck Arizona's order that a grower pack its premium cantaloupes in Arizona rather than in California, announcing the \textit{Pike} balancing test for evenhanded laws that burden interstate commerce (\href{https://supreme.justia.com/cases/federal/us/397/137/}{397 U.S. 137}).

\section{Detailed Case Facts}
Bruce Church grew superior cantaloupes at Parker, Arizona, and trucked them 31 miles to its efficient packing sheds in Blythe, California, where they were packed under equivalent standards in crates bearing a California name. Arizona's enforcement official, Loren Pike, ordered the melons packed in Arizona, which would have required building a \$200,000 plant and risked a \$700,000 crop.

\section{Procedural History}
A three-judge district court enjoined the order; Pike appealed directly, and the Supreme Court affirmed.

\section{Judicial Votes}
Affirmed, 8--0. Justice Stewart wrote for a unanimous Court, with one seat vacant; there were no separate opinions.

\section{Holding}
An evenhanded statute serving a legitimate local interest with only incidental interstate effects stands unless its burden is ``clearly excessive in relation to the putative local benefits.'' Arizona's minimal interest in reputation cannot justify forcing a company to build an unneeded in-state plant.

\section{Analysis of Opinions}
Stewart held that the melons were in interstate commerce because they were destined for immediate out-of-state shipment, and Arizona demanded that an out-of-state operation be relocated in-state, a mandate condemned in \textit{Toomer v. Witsell} and the \textit{Haydel} shrimp cases. The order served no health, safety, or anti-deception purpose; Arizona sought only ``reflected good will'' for other growers, an interest ``minimal, at best.'' Forcing operations in-state that could be done more efficiently elsewhere is ``virtually per se illegal,'' a ``straitjacket'' on interstate resources. In \textit{National Pork Producers Council v. Ross} (2023), the Court fractured over \textit{Pike}, with a majority recasting it as a tool to ``smoke out'' protectionism but six Justices preserving balancing.

\section{Examples: Future Applications}
\textbf{Same-Side:} requiring shrimp to be peeled in-state before export is invalid. \textbf{Opposite-Side:} an evenhanded ban on polystyrene produce packaging, with no local-processing mandate, is valid. \textbf{Fence-Sitter:} a sales ban on milk from hormone-treated cows that forces costly out-of-state supply chains.

\section{Critique}
Though famed for balancing, \textit{Pike} rested heavily on the ``virtually per se illegal'' rule against forced local processing. It offers no metric for weighing costs against noneconomic values and ignored an obvious alternative, origin stamps on California-packed crates.

\section{Key Quotations}
\begin{quote}
``Although the criteria for determining the validity of state statutes affecting interstate commerce have been variously stated, the general rule that emerges can be phrased as follows: where the statute regulates evenhandedly to effectuate a legitimate local public interest, and its effects on interstate commerce are only incidental, it will be upheld unless the burden imposed on such commerce is clearly excessive in relation to the putative local benefits. \textit{Huron Cement Co. v. Detroit}, 362 U. S. 440, 443. If a legitimate local purpose is found, then the question becomes one of degree. And the extent of the burden that will be tolerated will, of course, depend on the nature of the local interest involved, and on whether it could be promoted as well with a lesser impact on interstate activities.''
\end{quote}
\noindent --- Justice Stewart, for the Court, 397 U.S. at 142.

\end{document}
